% ============================================================
% Paper 6 - Positive Curvature of the Spectral Trace
%           at the Critical Line
%
% Author:   Ulrich Tehrani
% Date:     April 2026
% Version:  1.2
% Zenodo:   https://doi.org/10.5281/zenodo.19665790
% License:  CC BY 4.0
%
% Changelog
% ---------
% v1.0 (April 2026)
%   First published LaTeX version.
%
% v1.1 (April 2026)
%   - Notation and Conventions section added between Abstract
%     and Introduction.
%   - Bibliography: Paper 5 title corrected to match the Zenodo
%     record ("Spectral Trace Formula and Smoothed Zero Sums:
%     A Prime-Zero Duality Framework"). DOI unchanged.
%   - Abstract Structure paragraph: "self-contained" replaced by
%     "non-circular and proof-complete" to remove an ambiguity
%     flagged in external review. The corresponding phrase in
%     the Reader Contract and Open Problems II-IV is retained,
%     since the modifier "within classical analytic number
%     theory" disambiguates in those contexts.
%   - PDF Unicode hygiene: \usepackage{cmap} added before
%     \usepackage[T1]{fontenc}. This fixes the ToUnicode
%     mapping for the section sign (T1 slot 0x9F) and en/em-
%     dashes (T1 slots 0x15 and 0x16) in the PDF text stream;
%     visual rendering is unchanged.
%   - Source hygiene: internal project codename removed from
%     the Call for Feedback section ("the project" -> "the
%     series") and from the bibliography entry for the
%     verification code ("Verification scripts for the
%     series"). URLs remain unchanged.
%   - Motto placement: the motto "Pauca sed matura", which
%     previously appeared three times (after Acknowledgments,
%     after Call for Feedback, in the footer), now appears
%     exactly once, in the footer block. Used sparingly the
%     motto carries weight; repeated, it dilutes.
%   - No changes to mathematical content; no numerical values
%     changed.
%
% v1.2 (April 2026)
%   - Literature positioning strengthened based on two external
%     Deep Research passes (citation verification + positioning
%     & gap analysis).
%   - Edwards pinpoint verified (§1.16 and §3.6 already present
%     in v1.1; confirmed correct). Weil "Suppl." already present.
%   - Eight positioning references added: Connes (1999),
%     Berry-Keating (1999), de Branges (1986), Burnol (2004),
%     Bombieri (2000), Li (1997), Bombieri-Lagarias (1999),
%     Rubinstein-Sarnak (1994).
%   - Three positioning sentences added: operator-programme
%     distinction (Introduction), curvature-vs-Li comparison
%     (Main Results), scope-and-comparison (Open Problems).
%   - Prime-field narrative clarified with explicit dictionary
%     disclaimer.
%   - Re(Z̃_p) < 0 distinguished from Chebyshev bias.
%   - Call for Feedback updated with diagnostic expert question.
%   - Three precision fixes: Re(Z̃_p) = Re(Z_p) equality replaced
%     by truncation-controlled approximation; O(ε) notation in
%     Theorem 4.1 and Abstract replaced by numerical wording;
%     "converges" in Theorem 5.1(a) weakened to "appears to
%     converge (numerical extrapolation)".
%   - No changes to mathematical content, proofs, or numerical
%     values.
% ============================================================

\documentclass[11pt]{article}
\usepackage{cmap}
\usepackage[T1]{fontenc}
\usepackage[utf8]{inputenc}
\usepackage[a4paper,margin=1in]{geometry}
\usepackage{amsmath,amssymb,amsthm,mathtools}
\usepackage{xcolor}
\usepackage{hyperref}
\usepackage{enumitem}
\usepackage{setspace}
\usepackage{booktabs}
\usepackage{array}
\usepackage{graphicx}

\hypersetup{colorlinks=true,
  linkcolor=blue!50!black,
  urlcolor=blue!50!black,
  citecolor=blue!50!black}
\setlength{\parskip}{0.5em}
\setlength{\parindent}{0pt}
\onehalfspacing

% --------------------------------------------------------
% Theorem environments (numbered within section)
% --------------------------------------------------------
\newtheorem{theorem}{Theorem}[section]
\newtheorem{proposition}[theorem]{Proposition}
\newtheorem{lemma}[theorem]{Lemma}
\newtheorem{corollary}[theorem]{Corollary}
\newtheorem{definition}[theorem]{Definition}
\newtheorem{assumption}[theorem]{Assumption}
\newtheorem{observation}[theorem]{Observation}
\newtheorem{conjecture}[theorem]{Conjecture}
\newtheorem*{remark}{Remark}
\newtheorem{openproblem}[theorem]{Open Problem}

% Heuristic environment (clearly marked, never used in proofs)
\newenvironment{heuristic}[1]{%
  \par\smallskip\noindent\textit{[Heuristic: #1]}\quad}{%
  \par\smallskip}

% --------------------------------------------------------
% Convention box — clean ruled box, no mdframed
% --------------------------------------------------------
\newenvironment{conventionbox}{%
  \medskip
  \noindent\rule{\linewidth}{0.4pt}\smallskip\par\noindent
}{%
  \par\smallskip\noindent\rule{\linewidth}{0.4pt}
  \medskip
}

% --------------------------------------------------------
% Series-wide macros
% --------------------------------------------------------
\newcommand{\Hstr}{H_{\mathrm{str}}}
\newcommand{\Hnull}{H_{\mathrm{null}}}
\newcommand{\Estr}{E_{\mathrm{str}}}
\newcommand{\Espec}{E_{\mathrm{spec}}}
\newcommand{\etaorig}{\eta_{\mathrm{orig}}}
\newcommand{\Tren}{T_{\mathrm{ren}}}
\newcommand{\Tt}{\widetilde{T}}
\newcommand{\Gun}{G^{\mathrm{un}}}
\newcommand{\Hlocal}{H_{\mathrm{local}}}

% Smoothed zero sums and trace objects
\newcommand{\Zt}{\widetilde{Z}}
\newcommand{\DSEL}{D_{\mathrm{SEL}}}
\newcommand{\Ediag}{E_{\mathrm{diag}}}
\newcommand{\Ddiag}{D_{\mathrm{diag}}}

% Constants (always use subscript — never plain C)
\newcommand{\Ceta}{C_{\eta}}
\newcommand{\CT}{C_{T}}

% --------------------------------------------------------
% Title
% --------------------------------------------------------
\title{\vspace{-1.5em}\bfseries Positive Curvature of the Spectral Trace
at the Critical Line}
\author{Ulrich Tehrani}
\date{Research Note \textperiodcentered\ April 2026
      \textperiodcentered\ v1.2}

% ============================================================
\begin{document}
\maketitle
% ============================================================

\begin{abstract}
Building on the exact trace formula $\mathrm{Tr}(\Tt(\sigma))=\DSEL-O(\sigma)$
established in~\cite{Paper5}, we analyse the second-order curvature of the
spectral trace at the critical line.
The central quantity is the direct curvature sum
$B=\sum_{k,p} e^{-\varepsilon^2\gamma_k^2}(\gamma_k\log p)^2\cos(\gamma_k\log p)$,
which satisfies the algebraic identity $O''(\tfrac12)=-2B$, together with a
Guinand--Weil decomposition of a smoothed prime--zero sum $\Zt_p(\varepsilon)$
whose leading term is proved strictly negative.
No hypothetical input is used: no Riemann Hypothesis, no Hilbert--P\'olya
postulate, no GUE conjecture.

\textbf{What is proved.}
For every prime $p$ and every $\varepsilon>0$, the leading term of the
Guinand--Weil decomposition satisfies
$\mathrm{Main}_p(\varepsilon)<0$ (Theorem~\ref{thm:main-neg}), while the
other-prime error is non-positive
(Theorem~\ref{thm:other}) and the zero-sum truncation error at $N=100$
is bounded by $3\times 10^{-61}$ relative to the main term
(Theorem~\ref{thm:trunc}).
The algebraic identity $O''(\tfrac12)=-2B$ connecting curvature to bias
is derived from the trace formula of~\cite{Paper5}
(Proposition~\ref{prop:curv-bias}).
The decomposition
$\Zt_p(\varepsilon)=\mathrm{Main}_p+\Gamma_p+\mathrm{Const}_p
+\mathrm{Err}_{p^k}+\mathrm{Err}_{\mathrm{other}}$
is the classical Guinand--Weil explicit formula applied to an even,
Gaussian-damped test function; its derivation is recalled in~\S\,\ref{sec:gw}.

\textbf{What is numerical.}
The gamma factor contribution is subleading, with
$|\Gamma_p(\varepsilon)|/|\mathrm{Main}_p(\varepsilon)|$ numerically
consistent with linear scaling in~$\varepsilon$
(Theorem~\ref{thm:gamma}).
At the reference cutoff $\kappa=53$, the constant-term ratio satisfies
$r_p\in[0.75,0.80]$ for every prime $p\le 53$
(Theorem~\ref{thm:ratio-signcross}(a)).
Direct grid evaluation on the tested values
$\varepsilon\in\{0.020,\,0.025,\,0.030\}$ locates the sign-crossover
in the interval $(0.020,\,0.025)$, with
$\mathrm{Re}\,\Zt_p(\varepsilon)<0$ at the tested $\varepsilon=0.020$
for every prime $p\le 53$ (Theorem~\ref{thm:ratio-signcross}(b)).
The integrated bias is negative at reference parameters,
$B_{\mathrm{int}}(0.05)=-42.21<0$ at
$(\kappa,\varepsilon,N)=(53,0.05,100)$ (Theorem~\ref{thm:intbias}).
At the same reference parameters the direct curvature sum equals
$B=-19\,342.5<0$, which is equivalent to
$O''(\tfrac12)=+38\,685>0$ by the curvature--bias identity
(Corollary~\ref{cor:curvature}).

\textbf{What is conditional.}
If, in addition to the numerically established inequality $B<0$ at the
reference parameters, the stationarity condition $O'(\tfrac12)=0$
holds at the same parameters, then $\sigma=\tfrac12$ is a strict local
minimum of $O$ and $\mathrm{Tr}(\Tt(\sigma))$ attains a strict local
maximum at $\sigma=\tfrac12$ at the reference parameters
(Remark following Corollary~\ref{cor:curvature}).
The conditional statement is clearly marked; its analytic hypothesis is
the subject of Open Problem~\ref{op:stat}.

\textbf{What remains open.}
Five questions separate the present work from a full curvature proof,
each precisely named in~\S\,\ref{sec:open}.
Problem~\ref{op:stat} (\emph{Stationarity}) asks for an analytic bound
$|O'(\tfrac12)|\ll W\cdot P$, where $W=\sum_k e^{-\varepsilon^2\gamma_k^2}$
and $P=\pi(\kappa)$.
Problem~\ref{op:bias-asymp} (\emph{Asymptotic pointwise bias}) asks for
an analytic proof of $\mathrm{Re}\,\Zt_p(\varepsilon)<0$ for every prime
$p\le\kappa$.
Problem~\ref{op:uniformity} (\emph{Uniformity of positive curvature})
asks whether the inequalities $B<0$ and $O''(\tfrac12)>0$ persist at
parameter points $(\kappa,\varepsilon,N)$ other than the reference
values tested here.
Problem~\ref{op:transfer} (\emph{Integrated-to-direct transfer}) asks
whether pointwise or integrated bias implies the negativity of $B$.
Problem~\ref{op:weil-bridge} (\emph{Lagarias--Weil bridge}) asks whether
the curvature statement transfers to the Weil-positivity framework
of~\cite{Lagarias}.

\textbf{Structure.}
The formal core in \S\S\,\ref{sec:gw}--\ref{sec:curvcor} is non-circular
and proof-complete.
No use of the Riemann Hypothesis, the GUE conjecture, the Montgomery
pair correlation conjecture, or the Hilbert--P\'olya postulate is made
anywhere in \S\S\,\ref{sec:gw}--\ref{sec:curvcor}.
\end{abstract}

\vspace{2em}

% --------------------------------------------------------
% Notation and Conventions  (template invariant, v1.3)
% --------------------------------------------------------
\section*{Notation and Conventions}
\label{sec:notation}

\emph{Critical line.}
Throughout, the critical line means $\Re(s)=\tfrac12$.
The trace parameter $\sigma\in\mathbb{R}$ is evaluated at
$\sigma=\tfrac12$ unless stated otherwise.

\medskip
\emph{Global parameters.}
The reference values are $\kappa=53$ (prime cutoff, yielding the
$16$ active primes $p\le\kappa$), $\varepsilon=0.05$ (Gaussian
regularisation parameter), and $N=100$ (number of Riemann ordinates
$\gamma_k$ employed).
These values are fixed throughout
\S\S\,\ref{sec:gw}--\ref{sec:curvcor} unless explicitly varied.

\medskip
\emph{Main objects.}
The finite truncation
\[
Z_p \;=\; Z_{p,\varepsilon,N} \;:=\;
\sum_{k=1}^{N} e^{-\varepsilon^2\gamma_k^2}\,p^{i\gamma_k}
\]
is the numerical object of~\cite{Paper4,Paper5}.
The associated Guinand--Weil object
\[
\Zt_p(\varepsilon) \;:=\; \sum_{\rho}
h^{\mathrm{even}}_{p,\varepsilon}\!\left(\tfrac{\rho-1/2}{i}\right),
\qquad
h^{\mathrm{even}}_{p,\varepsilon}(u)=e^{-\varepsilon^2 u^2}\cos(u\log p),
\]
sums over all non-trivial zeros of $\zeta$.
Both quantities are real-valued by the symmetric pairing
$\rho\leftrightarrow\bar\rho$.
$\Re Z_{p,\varepsilon,N}$ approximates $\Re\Zt_p(\varepsilon)$ with
truncation tail $R_{p,N}(\varepsilon)$; at
$(\varepsilon,N)=(0.05,100)$ this tail is negligible relative to
$|\mathrm{Main}_p|$ by Theorem~\ref{thm:trunc}, hence all reported
numerics are reproducible via~$Z_{p,\varepsilon,N}$.
The direct curvature sum
\[
B \;:=\; \sum_{k,p} e^{-\varepsilon^2\gamma_k^2}\,
(\gamma_k\log p)^2\,\cos(\gamma_k\log p)
\]
enters the algebraic identity $O''(\tfrac12)=-2B$
(Proposition~\ref{prop:curv-bias}); the integrated bias is
$B_{\mathrm{int}}(\varepsilon)\;:=\;\sum_p (\log p)^2\,\Re\Zt_p(\varepsilon)$.

\medskip
\emph{Not to be confused with.}
$B$ (curvature sum, weight $(\gamma_k\log p)^2$) and $B_{\mathrm{int}}$
(integrated bias, weight $(\log p)^2$) differ in their weighting and
in the object summed; both are negative at the reference parameters,
but neither implies the other (Open Problem~\ref{op:transfer}).
$Z_p$ (finite truncation, level~$N$) and $\Zt_p$ (Guinand--Weil object,
sum over all non-trivial zeros) agree in real part for finite $N$ but
are formally distinct summations.
The diagonal-energy constants $D\approx 9.471$ of~\cite{Paper2} and
$\DSEL\approx 10.985$ of~\cite{Paper5} share a name across the series
but are defined in different operator contexts and must not be
identified.

\newpage

% -------------------------------------------------------
\section{Introduction}
\label{sec:intro}

\subsection{Background and Motivation}

This paper is the sixth in a series developing an operator-theoretic
framework for studying the distribution of the non-trivial zeros of
the Riemann zeta function $\zeta(s)$.
The series is built around a finite-dimensional Hilbert-space model in
which prime numbers mediate coupling between zeros, encoded in an
explicit operator family.

\textbf{Paper~1}~\cite{Paper1} introduced the local curvature function
$\Hlocal(\sigma,\kappa)=\sum_{p\le\kappa}V_p(\sigma)$, identifying
$\sigma=\tfrac12$ as a phase boundary through the divergence
$\Hlocal(\tfrac12,\kappa)\sim 8(\log\kappa)^2$.

\textbf{Paper~2}~\cite{Paper2} connected this local curvature to the
Weil explicit functional, constructing admissible test functions
for which the curvature divergence becomes a structural feature of
the Weil energy.

\textbf{Paper~3}~\cite{Paper3} made the geometric structure explicit:
two finite-dimensional Hilbert spaces $\Hstr$ (over primes) and
$\Hnull$ (over zeros) connected by a linear coupling map $\Phi$, with
loop operator $T=\Phi^*\circ\Phi$ and a first numerical proof of the
energy asymmetry $\etaorig>0$.

\textbf{Paper~4}~\cite{Paper4} introduced the dual operator
$\Tt=\Phi\circ\Phi^*$ on $\Hnull$ --- the Tehrani operator --- and
established a conditional proof of energy asymmetry under an
equidistribution assumption on the prime log-phases.

\textbf{Paper~5}~\cite{Paper5} extended this to a one-parameter
family $\Tt(\sigma)=\Phi(\sigma)\circ\Phi(\sigma)^*$ and proved the
exact spectral trace identity
\[
\mathrm{Tr}(\Tt(\sigma)) \;=\; \DSEL - O(\sigma),
\qquad
O(\sigma) \;=\; \tfrac12\sum_{k,p} e^{-\varepsilon^2\gamma_k^2}
                \cos(2\sigma\gamma_k\log p),
\]
where $\DSEL$ is a $\sigma$-independent constant.
Four independent numerical signatures concentrating at $\sigma=\tfrac12$
were documented in~\cite{Paper4,Paper5}, and the curvature-bias quantity
$B=\sum_p(\log p)^2\,\mathrm{Re}(Z_p)$ was introduced with the reference
value $B=-19\,342.5$ at $(\kappa,\varepsilon,N)=(53,0.05,100)$.

\textbf{The present paper} addresses the analytic question left open
by~\cite{Paper5}: \emph{is there a structural reason, independent of
the Riemann Hypothesis, why $\sigma=\tfrac12$ is distinguished by the
trace?}
We answer this in the affirmative at the level of second-order
curvature, at the reference parameters.
Twice-differentiating the trace identity at $\sigma=\tfrac12$ yields
the algebraic relation $O''(\tfrac12)=-2B$, so that the sign of $B$
determines the curvature of $O$ at the critical line.
At the reference parameters $B$ is numerically negative, hence
$O''(\tfrac12)$ is numerically positive.
The paper establishes this curvature statement together with a
Guinand--Weil decomposition of a related smoothed prime--zero sum,
provides four numerical witnesses of the same qualitative phenomenon,
and formulates five open problems constituting the completion path
toward a theorem-level selection principle.

The operator $\Tt(\sigma)$ should be viewed neither as a
Hilbert--P\'olya spectral realisation in the sense of
Connes~\cite{Connes} or Berry--Keating~\cite{BerryKeating}, nor as a
de~Branges structure-function criterion~\cite{deBranges}; it is a
finite-cutoff operator packaging Gaussian-smoothed explicit-formula
data, closest in spirit to Hilbert-space reformulations of the
explicit formula such as those of Burnol~\cite{Burnol}.

The heuristic labels ``prime field'' and ``field strength'' used in
the series denote a finite-cutoff coupling picture in which primes
mediate pairwise interactions between zero ordinates via explicit
Gaussian-smoothed phasors; this should be distinguished from the
periodic-orbit dictionary of Berry--Keating~\cite{BerryKeating},
the adelic trace-formula picture of Connes~\cite{Connes}, and the
Frobenius/cohomology picture of function-field zeta functions.

\subsection{Main Results}

The paper contains three main components.

The first is the curvature statement at the critical line.
At the reference parameters $(\kappa,\varepsilon,N)=(53,0.05,100)$,
direct numerical evaluation gives $B=-19\,342.5<0$, which is equivalent
via the algebraic identity of Proposition~\ref{prop:curv-bias} to
$O''(\tfrac12)=-2B=+38\,685>0$
(Corollary~\ref{cor:curvature}).
Independently, when combined with the stationarity hypothesis
$O'(\tfrac12)=0$, this inequality implies that $\sigma=\tfrac12$ is a
strict local minimum of $O$ at the reference parameters; the
stationarity hypothesis is stated explicitly and left as
Problem~\ref{op:stat} of~\S\,\ref{sec:open}.

The curvature identity $O''(\tfrac12)=-2B$ is structurally analogous
to the Li/Bombieri--Lagarias sign
criteria~\cite{Li,BombieriLagarias} in that it turns
explicit-formula data into a sign condition, but it is not a
derivative of $\xi$ at $s=1$ and not a Weil-positivity theorem; it
is a local second-order statement in the auxiliary trace parameter
$\sigma$ at the symmetry point~$\tfrac12$.

The second is the Guinand--Weil bias analysis of~\S\S\,\ref{sec:gw}--\ref{sec:const}.
The smoothed prime--zero sum $\Zt_p(\varepsilon)$ is decomposed into
five contributions, of which the leading one is proved negative
(Theorem~\ref{thm:main-neg}) and two further terms are proved
non-positive or exponentially small
(Theorems~\ref{thm:other},~\ref{thm:trunc}).
The gamma contribution is shown numerically to be subleading of order
$\varepsilon$ (Theorem~\ref{thm:gamma}).
The constant-term ratio $r_p=\lim_{\varepsilon\to 0}
|\mathrm{Const}_p(\varepsilon)|/|\mathrm{Main}_p(\varepsilon)|$
is numerically established to lie in $[0.75,0.80]$ for every prime
$p\le 53$ (Theorem~\ref{thm:ratio-signcross}(a)), and a sign-crossover
is localised in $(0.020,\,0.025)$
(Theorem~\ref{thm:ratio-signcross}(b)).
The integrated bias
$B_{\mathrm{int}}(\varepsilon)=\sum_{p\le\kappa}(\log p)^2\,
\mathrm{Re}\,\Zt_p(\varepsilon)$ is numerically negative at the
reference parameters, with $B_{\mathrm{int}}(0.05)=-42.21$
(Theorem~\ref{thm:intbias}).

The third is the structural separation made explicit in~\S\,\ref{sec:intbias}.
The quantities $B$, $B_{\mathrm{int}}(\varepsilon)$, and the pointwise
sign of $\mathrm{Re}\,\Zt_p(\varepsilon)$ are structurally distinct
witnesses of the same qualitative phenomenon: $B$ carries the
$(\gamma_k\log p)^2$-weighting, $B_{\mathrm{int}}$ only the
$(\log p)^2$-weighting, and the pointwise statement carries no
weighting at all.
All three are numerically negative at the reference parameters; the
analytical relationship between them is left open as
Problem~\ref{op:transfer}.
The five open problems of~\S\,\ref{sec:open} list the steps remaining,
in order of what we consider tractability:
Problem~\ref{op:stat} (stationarity) is the most tractable and, together
with the numerically established $B<0$, would suffice to upgrade
Corollary~\ref{cor:curvature} to an unconditional strict local-minimum
statement at the reference parameters.

\subsection{Reader Contract}

This paper \emph{proves}, by direct computation and classical estimates
of the Guinand--Weil type, that the main term is negative
(Theorem~\ref{thm:main-neg}), that the other-prime error is
non-positive (Theorem~\ref{thm:other}), that the truncation error is
negligible (Theorem~\ref{thm:trunc}), and the algebraic
curvature--bias identity $O''(\tfrac12)=-2B$
(Proposition~\ref{prop:curv-bias}).
These are unconditional mathematical statements.

This paper \emph{establishes numerically}, at the reference parameters
$(\kappa,\varepsilon,N)=(53,0.05,100)$, that the gamma contribution is
subleading of order $\varepsilon$ (Theorem~\ref{thm:gamma}), that
$r_p\in[0.75,0.80]$ for every prime $p\le 53$ and that the pointwise
bias $\mathrm{Re}\,\Zt_p<0$ holds on the tested subgrid at
$\varepsilon=0.020$ (Theorem~\ref{thm:ratio-signcross}), that
$B_{\mathrm{int}}(0.05)=-42.21<0$ (Theorem~\ref{thm:intbias}), and that
the positive-curvature statement $O''(\tfrac12)=+38\,685>0$ holds at
the reference parameters (Corollary~\ref{cor:curvature}).
These statements are verified on a finite grid and are clearly labelled
[numerical] in their theorem headings.

This paper \emph{states conditionally} that $\sigma=\tfrac12$ is a
strict local minimum of $O$ at the reference parameters, under the
hypothesis $O'(\tfrac12)=0$ (Remark following
Corollary~\ref{cor:curvature}).
The hypothesis is the object of Problem~\ref{op:stat}, and no part of
the unconditional claims above depends on it.

This paper \emph{leaves open} five well-posed questions named in
\S\,\ref{sec:open}: Problem~\ref{op:stat} (Stationarity),
Problem~\ref{op:bias-asymp} (Asymptotic pointwise bias),
Problem~\ref{op:uniformity} (Uniformity of positive curvature),
Problem~\ref{op:transfer} (Integrated-to-direct transfer), and
Problem~\ref{op:weil-bridge} (Lagarias--Weil bridge).
The first four are self-contained within classical analytic number
theory; the fifth names the remaining bridge to Weil positivity in
the formulation of~\cite{Lagarias}.

% -------------------------------------------------------
\section{The Guinand--Weil Decomposition}
\label{sec:gw}

Throughout, $p$ denotes a rational prime, $\varepsilon>0$ a smoothing
parameter, and $\rho=\tfrac12+i\gamma$ a non-trivial zero of the
Riemann zeta function.
We fix a cutoff $\kappa\ge 2$ and consider only primes $p\le\kappa$.
For numerical reference we use $(\kappa,\varepsilon,N)=(53,0.05,100)$,
where $N$ is the number of zeros used in the finite approximation;
the first $16$ primes lie below $\kappa=53$.

\begin{definition}[Smoothed zero sum]
\label{def:smoothed}
Let
\[
Z_p \;=\; Z_{p,\varepsilon,N} \;:=\; \sum_{k=1}^{N} e^{-\varepsilon^2\gamma_k^2}\,
p^{i\gamma_k},
\qquad
\mathrm{Re}\,Z_p \;=\; \sum_{k=1}^{N} e^{-\varepsilon^2\gamma_k^2}
                        \cos(\gamma_k\log p).
\]
This is the finite numerical object studied in~\cite{Paper4,Paper5}.
The associated object obtained from the Guinand--Weil explicit formula is
\[
\Zt_p(\varepsilon) \;:=\; \sum_{\rho} h^{\mathrm{even}}_{p,\varepsilon}
\!\left(\tfrac{\rho-\tfrac12}{i}\right),
\qquad
h^{\mathrm{even}}_{p,\varepsilon}(u) \;=\; e^{-\varepsilon^2 u^2}
                                            \cos(u\log p).
\]
\end{definition}

Since $h^{\mathrm{even}}_{p,\varepsilon}$ is even and decays rapidly,
$\Zt_p(\varepsilon)$ is real-valued and absolutely convergent.
Its Fourier transform is
\[
\widehat{h}_{p,\varepsilon}(x) \;=\; \frac{\sqrt{\pi}}{2\varepsilon}
\!\left[\exp\!\left(-\tfrac{(\log p-2\pi x)^2}{4\varepsilon^2}\right)
      +\exp\!\left(-\tfrac{(\log p+2\pi x)^2}{4\varepsilon^2}\right)
\right],
\]
which is strictly positive on $\mathbb R$ and concentrates at
$\pm\log p/(2\pi)$ as $\varepsilon\to 0$.

\begin{proposition}[Guinand--Weil explicit formula]
\label{prop:gw}
For every prime $p$ and every $\varepsilon>0$,
\[
\Zt_p(\varepsilon) \;=\; \mathrm{Main}_p(\varepsilon)
                       + \Gamma_p(\varepsilon)
                       + \mathrm{Const}_p(\varepsilon)
                       + \mathrm{Err}_{p^k}(\varepsilon)
                       + \mathrm{Err}_{\mathrm{other}}(\varepsilon),
\]
where the five contributions are defined by the standard partition of
the prime sum and the archimedean terms of the explicit formula.
Explicitly,
\[
\mathrm{Main}_p(\varepsilon) \;=\; -\frac{1}{2\pi}\cdot\frac{\log p}{\sqrt p}
\!\left[\widehat{h}_{p,\varepsilon}\!\left(\tfrac{\log p}{2\pi}\right)
      +\widehat{h}_{p,\varepsilon}\!\left(-\tfrac{\log p}{2\pi}\right)
\right],
\]
\[
\Gamma_p(\varepsilon) \;=\; \frac{1}{2\pi}\int_{-\infty}^{\infty}
e^{-\varepsilon^2 t^2}\cos(t\log p)\,
\mathrm{Re}\,\psi\!\left(\tfrac14 + \tfrac{it}{2}\right)\,dt,
\]
\[
\mathrm{Const}_p(\varepsilon) \;=\; \Zt_p(\varepsilon)
\,-\, \mathrm{Main}_p(\varepsilon)
\,-\, \Gamma_p(\varepsilon)
\,-\, \mathrm{Err}_{p^k}(\varepsilon)
\,-\, \mathrm{Err}_{\mathrm{other}}(\varepsilon),
\]
where $\psi$ denotes the digamma function.
The prime-power error collects the contribution from $n=p^k$ with
$k\ge 2$, and the remaining-prime error from $n=q^k$ with $q\neq p$;
both are standard sums of the form
$\sum_n\Lambda(n)\widehat{h}_{p,\varepsilon}(\log n/(2\pi))/\sqrt n$.
\end{proposition}

\begin{remark}
At the poles $s=0,1$ of $\xi$, the test function contributes the value
$h^{\mathrm{even}}_{p,\varepsilon}(u)$ evaluated at $u=\pm i/2$.
A direct computation gives
\[
h^{\mathrm{even}}_{p,\varepsilon}\!\left(\tfrac{i}{2}\right)
+ h^{\mathrm{even}}_{p,\varepsilon}\!\left(-\tfrac{i}{2}\right)
\;=\; 2\,e^{\varepsilon^2/4}\cosh\!\left(\tfrac{\log p}{2}\right)
\;=\; e^{\varepsilon^2/4}\!\left(\sqrt p + \tfrac{1}{\sqrt p}\right).
\]
This is the hyperbolic cosine, not the circular cosine: the formal
substitution $\cos(i\log p/2)=\cosh(\log p/2)$ is essential.
Replacing $\cosh$ by $\cos$ would yield a quantity that is negative
for $p\ge 29$, which is analytically impossible since the left-hand
side equals $\sqrt p + 1/\sqrt p>2$ for every $p\ge 2$.
The correct $\cosh$ form is used throughout what follows.
See Edwards~\cite[\S1.16 and \S3.6]{Edwards} for the classical derivation
of the pole contributions at $s=0,1$.
\end{remark}

The remainder of the paper studies the five terms of the decomposition
in turn.

% -------------------------------------------------------
\section{The Main Term}
\label{sec:main}

\begin{theorem}[Main term negativity; proved]
\label{thm:main-neg}
For every prime $p$ and every $\varepsilon>0$,
\[
\mathrm{Main}_p(\varepsilon) \;=\; -\frac{\log p}{2\sqrt\pi\,\varepsilon\,\sqrt p}
\cdot\bigl(1 + O(e^{-(\log p)^2/\varepsilon^2})\bigr)
\;<\; 0.
\]
\end{theorem}

\begin{proof}
The Fourier transform $\widehat{h}_{p,\varepsilon}(x)$ is strictly
positive on $\mathbb R$ and attains its maximum at
$x=\pm\log p/(2\pi)$ with value $\sqrt\pi/\varepsilon$.
The reflected value at $x=\mp\log p/(2\pi)$ is smaller by the factor
$\exp(-(\log p)^2/\varepsilon^2)$, which is exponentially small: at
$\varepsilon=0.05$ and $p=2$ it is below $10^{-82}$, and decreases
further for larger $p$.
Collecting the dominant contribution,
\[
\mathrm{Main}_p(\varepsilon) \;=\; -\frac{1}{2\pi}\cdot\frac{\log p}{\sqrt p}
\cdot\frac{\sqrt\pi}{\varepsilon}
\bigl(1 + O(e^{-(\log p)^2/\varepsilon^2})\bigr)
\;=\; -\frac{\log p}{2\sqrt\pi\,\varepsilon\,\sqrt p}\bigl(1+o(1)\bigr).
\]
The sign is negative because all factors in the leading expression are
positive and enter with a minus sign.
\end{proof}

\begin{corollary}
\label{cor:main-scaling}
As $\varepsilon\to 0$, $|\mathrm{Main}_p(\varepsilon)|\sim
(\log p)/(\varepsilon\sqrt p)$, giving a leading-order scaling of
$1/\varepsilon$.
The overall constant is $C_{\mathrm{Main}}=-1/(2\sqrt\pi)\approx-0.2821$.
\end{corollary}

At the reference parameter $\varepsilon=0.05$ and for $2\le p\le 53$,
the main term takes values in the interval $[-4.15,-2.77]$, minimised
at $p=7$ and maximised at $p=2$.

% -------------------------------------------------------
\section{Error Terms}
\label{sec:errors}

We now treat the three remaining contributions other than the constant
term.

\subsection{The gamma contribution}

\begin{theorem}[Gamma term is subleading; numerical]
\label{thm:gamma}
For every prime $p\le\kappa$ and every sufficiently small
$\varepsilon>0$, the ratio
$|\Gamma_p(\varepsilon)|/|\mathrm{Main}_p(\varepsilon)|$
is numerically consistent with linear scaling in $\varepsilon$ on the
tested grid; a log--log fit over
$\varepsilon\in[0.005,\,0.1]$ yields slope~$1.001$.
At $\varepsilon=0.05$, the ratio is bounded above by $0.172$ (attained
at $p=2$) and below by $0.017$ (attained at $p=53$); in particular
$|\Gamma_p(0.05)|<|\mathrm{Main}_p(0.05)|$ for all $16$ primes $p\le 53$.
\end{theorem}

\begin{proof}[Sketch]
The integrand in $\Gamma_p(\varepsilon)$ is uniformly integrable on
$\mathbb R$ by the classical asymptotic
$\mathrm{Re}\,\psi(\tfrac14+it/2)=\log(1+|t|/2)+O(1)$ for large $|t|$
and the Gaussian weight $e^{-\varepsilon^2 t^2}$.
Dominated convergence gives $\Gamma_p(\varepsilon)\to\Gamma_p(0)$ as
$\varepsilon\to 0$, a finite limit.
Since $|\mathrm{Main}_p(\varepsilon)|$ grows as $1/\varepsilon$, the
ratio decays linearly.
A log--log regression of the ratio against $\varepsilon$ across
$\varepsilon\in[0.005,0.1]$ gives a slope of $1.001$, consistent with
the linear prediction to four significant figures.
The analytic step that would upgrade this to a proved bound is the
explicit control of $\Gamma_p(0)$ by a uniform constant in $p$, which
we do not establish here; it is documented as part of
Problem~\ref{op:bias-asymp} in~\S\,\ref{sec:open}.
\end{proof}

\subsection{The prime-power and other-prime errors}

\begin{theorem}[Other-prime error is non-positive; proved]
\label{thm:other}
\[
\mathrm{Err}_{\mathrm{other}}(\varepsilon) \;\le\; 0 \qquad
\text{for every } \varepsilon>0.
\]
\end{theorem}

\begin{proof}
By the standard form of the explicit formula,
\[
\mathrm{Err}_{\mathrm{other}}(\varepsilon) \;=\; -\frac{1}{2\pi}
\sum_{\substack{q\le\kappa\\ q\neq p\\ q\text{ prime}}}
\frac{\log q}{\sqrt q}
\!\left[\widehat{h}_{p,\varepsilon}\!\left(\tfrac{\log q}{2\pi}\right)
      +\widehat{h}_{p,\varepsilon}\!\left(-\tfrac{\log q}{2\pi}\right)
\right]
+ (\text{higher-prime-power terms}).
\]
Each summand has $\log q>0$, $\sqrt q>0$, and
$\widehat{h}_{p,\varepsilon}(\cdot)>0$ on $\mathbb R$, with an overall
minus sign.
The higher-prime-power terms share the same sign structure with
$\Lambda(q^k)=\log q$ and the same positive transform.
Hence the full sum is non-positive.
\end{proof}

Theorem~\ref{thm:other} is a structural identity rather than an error
bound: the contribution from primes $q\neq p$ actively reinforces the
negativity of $\Zt_p$.

\begin{theorem}[Truncation error is negligible; proved]
\label{thm:trunc}
\[
\frac{|R_{p,N}(\varepsilon)|}{|\mathrm{Main}_p(\varepsilon)|}
\;\le\; \exp(-\varepsilon^2\gamma_N^2)\cdot
       \frac{2\sqrt\pi\,\varepsilon\,\sqrt p}{\log p}
\;\le\; 3\times 10^{-61}
\]
at the reference parameters $(\varepsilon,N)=(0.05,100)$ for every
prime $p\le 53$, where $R_{p,N}(\varepsilon)=\sum_{k>N}
e^{-\varepsilon^2\gamma_k^2}\cos(\gamma_k\log p)$ is the tail of the
zero sum.
\end{theorem}

\begin{proof}
The Gaussian weight $e^{-\varepsilon^2\gamma_k^2}$ majorises the tail.
With $\gamma_{100}=236.524$ and $\varepsilon=0.05$ one has
$\varepsilon^2\gamma_{100}^2=139.77$, so that $e^{-139.77}<1.8\times 10^{-61}$.
The bound follows from the triangle inequality and the explicit lower
bound on $|\mathrm{Main}_p|$ from~\S\,\ref{sec:main}.
\end{proof}

For every downstream claim in this paper, truncation at $N=100$ is
numerically equivalent to summing the full zero series.

% -------------------------------------------------------
\section{The Constant Term and the Sign-Crossover Localisation}
\label{sec:const}

The constant term $\mathrm{Const}_p(\varepsilon)$ collects the pole
residues of $\xi$, the $\log(2\pi)$ normalisation, and the finite
remainder of the digamma contribution that is not captured in
$\Gamma_p$.
It is the only term of the decomposition that resists a short analytic
treatment.
We state a numerical result on its asymptotic ratio to the main term
and, independently, a numerical result localising the sign-crossover
of $\mathrm{Re}\,\Zt_p$ on a finite $\varepsilon$-grid.

\begin{theorem}[Asymptotic ratio bound and sign-crossover localisation; numerical]
\label{thm:ratio-signcross}
Numerically, the following hold.

\emph{(a)} For every prime $p\le 53$, the ratio
$|\mathrm{Const}_p(\varepsilon)|/|\mathrm{Main}_p(\varepsilon)|$
appears to converge (numerical extrapolation) to a finite limit
\[
r_p \;:=\; \lim_{\varepsilon\to 0}
\frac{|\mathrm{Const}_p(\varepsilon)|}{|\mathrm{Main}_p(\varepsilon)|}
\]
as $\varepsilon\to 0$, with $r_p\in[0.75,\,0.80]$.

\emph{(b)} Direct evaluation on the tested values
$\varepsilon\in\{0.020,\,0.025,\,0.030\}$ locates the sign-crossover
in the interval $(0.020,\,0.025)$.
At the tested value $\varepsilon=0.020$, one finds
\[
\mathrm{Re}\,\Zt_p(\varepsilon) \;<\; 0 \qquad
\text{for every prime } p\le 53.
\]

Both statements are numerical and refer to the cutoff $\kappa=53$.
The location of the sign-crossover and the sign at smaller tested
values of $\varepsilon$ are empirical findings on a finite grid; no
continuous interval-persistence statement is asserted.
The analytical question behind (a) is documented as
Problem~\ref{op:bias-asymp} of~\S\,\ref{sec:open}.
\end{theorem}

\emph{Discussion.}
The leading-order analysis of the pole-residue contribution at $s=1$
of $\xi$ uses the $\cosh$-evaluation of the remark in~\S\,\ref{sec:gw}
and yields a formal expression of order $\varepsilon^{-1}$, matching
the scaling of $\mathrm{Main}_p$ (see
Edwards~\cite[\S1.16 and \S3.6]{Edwards}).
The cancellation that determines the coefficient ratio is between
$e^{\varepsilon^2/4}(\sqrt p + 1/\sqrt p)$ on the pole side and
$-\log p/(2\sqrt\pi\,\varepsilon\,\sqrt p)$ on the main-term side.
Numerical evaluation of $\mathrm{Const}_p(\varepsilon)$ for
$\varepsilon\in\{0.005,0.010,0.015,0.020,0.025,0.030,0.050\}$ across
all primes $p\le 53$ shows that the ratio
$|\mathrm{Const}_p(\varepsilon)|/|\mathrm{Main}_p(\varepsilon)|$ is a
unimodal function of $\varepsilon$, with a maximum near
$\varepsilon\approx 0.045$--$0.060$ and limiting value
$r_p\in[0.75,0.80]$ as $\varepsilon\to 0$.
For $p\in\{37,53\}$ this maximum narrowly exceeds unity (attaining
$1.183$ and $1.216$ respectively), which produces two pointwise
exceptions to the bias at $\varepsilon=0.05$; see below.
The sign-crossover localisation in (b) is obtained from the grid
computation directly, not as an analytic consequence of (a).

\emph{Numerical evidence.}
At $\varepsilon=0.05$ and $\kappa=53$, the constant terms satisfy
$\mathrm{Const}_p\in[2.44,3.93]$ with mean $3.27$ and standard
deviation $0.46$ (coefficient of variation $14\%$).
Representative values are listed below.

\begin{table}[h]
\centering
\begin{tabular}{ccccc}
\toprule
$p$ & $\mathrm{Main}_p(0.05)$ & $\Gamma_p(0.05)$
    & $\mathrm{Const}_p(0.05)$ & $r_p(\varepsilon\to 0)$ \\
\midrule
$2$  & $-2.77$ & $-0.475$ & $2.44$ & $0.766$ \\
$7$  & $-4.15$ & $-0.193$ & $3.47$ & $0.759$ \\
$23$ & $-3.69$ & $-0.105$ & $3.79$ & $0.754$ \\
$37$ & $-3.35$ & $-0.082$ & $3.93$ & $0.796$ \\
$53$ & $-3.08$ & $-0.069$ & $3.74$ & $0.753$ \\
\bottomrule
\end{tabular}
\end{table}

At $\varepsilon=0.05$, the pointwise inequality
$\mathrm{Re}\,\Zt_p(0.05)<0$ holds for $14$ of the $16$ primes
$p\le 53$.
The exceptions are $p=37$ and $p=53$, for which
$\mathrm{Re}\,\Zt_p(0.05)\approx+0.50$ and $+0.59$ respectively.
This is \emph{not} a failure of the underlying decomposition: it
reflects the finite-$\varepsilon$ maximum of the ratio $r(\varepsilon)$
near $\varepsilon\approx 0.05$, which narrowly exceeds unity for
exactly these two primes.
At the tested value $\varepsilon=0.020$ the pointwise bias holds for
all $16$ primes without exception.
At $\varepsilon=0.025$ and $\varepsilon=0.030$ the exceptions
$p=37,53$ are present.
The sign-crossover thus lies in the interval $(0.020,\,0.025)$ as
established by these three tested values.
Whether the bias persists continuously throughout
$\varepsilon\in(0,0.020]$ is not established here.

No structural explanation for the exceptions $p=37,53$ at
$\varepsilon=0.05$ is currently available.
The phenomenon appears to be a finite-$\varepsilon$, finite-$\kappa$
manifestation of the observed unimodal peak of $r(\varepsilon)$ near
$\varepsilon\approx 0.045$--$0.060$, which for exactly these two
primes narrowly exceeds unity.
Whether this reflects an underlying arithmetic structure or is a
numerical coincidence at $\kappa=53$ is not addressed here.

The observed negativity of $\mathrm{Re}\,\Zt_p(\varepsilon)$ should
be read neither as a prime race in the sense of
Rubinstein--Sarnak~\cite{RubinsteinSarnak} nor as a Landau--Gonek
prime-power asymptotic; it is a different, prime-indexed sign
statistic for Gaussian-smoothed explicit-formula evaluations.

\emph{Status.}
The statement $r_p<1$ is verified numerically.
A formal proof --- which would close the asymptotic bias conjecture in
its sharpest form --- requires an explicit control of the
residue-series expansion near the poles $s=0,1$ and is part of
Problem~\ref{op:bias-asymp} of~\S\,\ref{sec:open}.
The empirical sign-crossover localisation above, at $\kappa=53$, is
the sharpest statement that the present methods allow.

% -------------------------------------------------------
\section{The Integrated Bias}
\label{sec:intbias}

The two expressions for $B$ used below are the same quantity in two
representations: the primal sum
\[
B \;=\; \sum_{k,p} e^{-\varepsilon^2\gamma_k^2}\,(\gamma_k\log p)^2\,
         \cos(\gamma_k\log p)
\]
and the dual sum $B=\sum_p(\log p)^2\,\mathrm{Re}\,Z_p$.
The equivalence between the two is Proposition~5.3 of~\cite{Paper5}
and is treated here as imported.
All numerical values in this section refer to the same object.

We use the dual representation as a starting point.
The curvature bias at $\sigma=\tfrac12$ is encoded in the weighted sum
\[
B \;:=\; \sum_{p\le\kappa}(\log p)^2\,\mathrm{Re}\,Z_p,
\]
which was introduced in~\cite{Paper5} as $-\tfrac12$ times the
$\sigma=\tfrac12$ evaluation of the second $\sigma$-derivative of the
trace identity, via the identity $O''(\tfrac12)=-2B$.
Its structural decomposition is algebraic and does not depend on any
unproved hypothesis.

\begin{proposition}[Curvature--bias identity; proved]
\label{prop:curv-bias}
Let $O(\sigma)=\tfrac12\sum_{k,p} e^{-\varepsilon^2\gamma_k^2}
\cos(2\sigma\gamma_k\log p)$ and define
\[
B \;:=\; \sum_{k,p} e^{-\varepsilon^2\gamma_k^2}\,(\gamma_k\log p)^2\,
         \cos(\gamma_k\log p).
\]
Then
\[
O''(\tfrac12) \;=\; -2\,B.
\]
\end{proposition}

\begin{proof}
Direct differentiation gives
\[
O'(\sigma) \;=\; -\sum_{k,p} e^{-\varepsilon^2\gamma_k^2}\,\gamma_k\log p\,
                \sin(2\sigma\gamma_k\log p),
\]
\[
O''(\sigma) \;=\; -2\sum_{k,p} e^{-\varepsilon^2\gamma_k^2}\,
                (\gamma_k\log p)^2\,\cos(2\sigma\gamma_k\log p).
\]
Evaluation at $\sigma=\tfrac12$ yields
\[
O''(\tfrac12) \;=\; -2\sum_{k,p} e^{-\varepsilon^2\gamma_k^2}\,
                (\gamma_k\log p)^2\,\cos(\gamma_k\log p) \;=\; -2\,B.
\qedhere
\]
\end{proof}

Numerically, at reference parameters $(\kappa,\varepsilon,N)=(53,0.05,100)$,
the direct curvature sum equals $B=-19\,342.5$, hence
$O''(\tfrac12)=-2B=+38\,685$.
This value coincides with $\mathrm{Tr}(L)$, where
$L=-d^2\Tt(\sigma)/d\sigma^2\big|_{\sigma=1/2}$ is the curvature
operator of~\cite{Paper5}, and provides an independent closed-form
cross-check for the identity above.

\begin{remark}
The quantity $B$ is defined as in~\cite[\S5]{Paper5}, and the identity
$O''(\tfrac12)=-2B$ follows from the explicit form of $O(\sigma)$
given there.
The numerical value $B=-19\,342.5$ at $(\kappa,\varepsilon,N)=(53,0.05,100)$
is the one quoted throughout the series.
\end{remark}

\begin{theorem}[Integrated bias is negative at reference parameters; numerical]
\label{thm:intbias}
At $(\kappa,\varepsilon,N)=(53,0.05,100)$,
\[
B_{\mathrm{int}}(0.05) \;:=\; \sum_{p\le 53}(\log p)^2\,
\mathrm{Re}\,\Zt_p(0.05) \;=\; -42.21 \;<\; 0.
\]
The direct curvature sum $B$ of Proposition~\ref{prop:curv-bias} is
likewise strictly negative at the same parameters, with
$B=-19\,342.5$, hence $O''(\tfrac12)=-2B=+38\,685$.
\end{theorem}

\begin{remark}
The values $B$ and $B_{\mathrm{int}}$ refer to two \emph{different}
quantities, not to a single quantity in two forms.
$B$ is the direct curvature sum weighted by $(\gamma_k\log p)^2$,
obtained by differentiating the trace identity and evaluating at
$\sigma=\tfrac12$; $B_{\mathrm{int}}$ is the integrated Guinand--Weil
bias, weighted only by $(\log p)^2$ and summed over
$\mathrm{Re}\,\Zt_p$.
The two differ by a reweighting through the $\gamma_k^2$-factors
absent in the second.
Both are numerically negative at the reference parameters, and the
shared negative sign is a numerical finding rather than the consequence
of a proved equivalence.
Proposition~\ref{prop:curv-bias} connects $B$ --- not $B_{\mathrm{int}}$
--- to the curvature $O''(\tfrac12)$.
Whether pointwise negativity of $\mathrm{Re}\,\Zt_p(\varepsilon)$, or
the integrated bias $B_{\mathrm{int}}(\varepsilon)<0$, implies the
corresponding negativity of $B$ at the same parameters is an open
question; it is formulated as Problem~\ref{op:transfer} of~\S\,\ref{sec:open}.
\end{remark}

\begin{remark}
Although $\mathrm{Re}\,\Zt_p(0.05)>0$ at $p=37$ and $p=53$, the
$(\log p)^2$-weighted sum $B_{\mathrm{int}}$ remains decisively
negative: the negative contributions from the remaining $14$ primes
dominate the positive contributions from these two.
Concretely, $(\log 37)^2\cdot 0.50+(\log 53)^2\cdot 0.59\approx
6.52+9.30=15.82$, against a cumulative negative contribution of
$-58.03$ from the other primes, yielding the net value $-42.21$.
The integrated bias is robust against individual-prime deviations as
long as these remain isolated.
\end{remark}

% -------------------------------------------------------
\section{Positive Curvature at the Critical Line}
\label{sec:curvcor}

We now record the curvature statement that follows from
Theorem~\ref{thm:intbias} and Proposition~\ref{prop:curv-bias}, and,
separately, the conditional strict-local-extremum statement that
follows once stationarity is added as a hypothesis.

\begin{corollary}[Positive curvature at $\sigma=\tfrac12$; numerical]
\label{cor:curvature}
At the reference parameters $(\kappa,\varepsilon,N)=(53,0.05,100)$,
the direct curvature sum $B$ is numerically negative
(Theorem~\ref{thm:intbias}), $B=-19\,342.5<0$.
By Proposition~\ref{prop:curv-bias}, this is equivalent to
\[
O''(\tfrac12) \;=\; -2B \;=\; +38\,685 \;>\; 0.
\]
\end{corollary}

\begin{proof}
By Theorem~\ref{thm:intbias}, $B=-19\,342.5$.
By Proposition~\ref{prop:curv-bias}, $O''(\tfrac12)=-2B$.
\end{proof}

\begin{remark}
If, in addition to $B<0$ at the reference parameters (here numerically
established), the stationarity condition $O'(\tfrac12)=0$ holds at the
same parameters, then $\sigma=\tfrac12$ is a strict local minimum of
$O$, and the trace
\[
\mathrm{Tr}\,\Tt(\sigma) \;=\; \DSEL - O(\sigma)
\]
attains a strict local maximum at $\sigma=\tfrac12$.

The vanishing of $O'(\tfrac12)$ is left open and formulated as
Problem~\ref{op:stat} of~\S\,\ref{sec:open}.
Numerically, $O'(\tfrac12)=-2.475$ is approximately $11\%$ of the
natural scale $W\cdot P$, consistent with the cancellation expected
from the odd kernel $\sin(\gamma_k\log p)$, but no analytical bound on
$|O'(\tfrac12)|$ is established in the present paper.
\end{remark}

\begin{remark}
The paper~\cite{Paper5} recorded four further numerical signatures
concentrating at $\sigma=\tfrac12$: a $+47\%$ trace spike relative to
the $\sigma$-neighbourhood average, the dominance of the leading
eigenvalue $\mu_1$, the minimum of the spectral gap, and the joint
maximum of $\mathrm{Tr}(L)$ and $\lambda_{\max}(L)$.
Each of these is numerically robust and independent of the bias
conditions established here.
They constitute additional numerical signatures from~\cite{Paper5} at
the same reference parameters, but do not enter the proof of
Corollary~\ref{cor:curvature} or of the conditional remark above.
\end{remark}

\begin{remark}
The curvature statement identifies $\sigma=\tfrac12$ as a geometrically
distinguished point of the spectral trace at the reference parameters.
It is a structural ingredient in a possible route to Weil positivity
in the Lagarias formulation~\cite{Lagarias}
(Problem~\ref{op:weil-bridge} of~\S\,\ref{sec:open}), but does not by
itself establish it.
No implication concerning the Riemann Hypothesis is asserted or
implied by Corollary~\ref{cor:curvature}, nor by the conditional remark
above when the stationarity hypothesis is imposed.
\end{remark}

% -------------------------------------------------------
\section{Open Problems}
\label{sec:open}

The present work reduces the question of curvature at the critical
line to five well-posed subproblems.
We name each and describe the tools we expect to be relevant.
The problems are listed in order of what we consider tractability.

\begin{openproblem}[Stationarity]
\label{op:stat}
Prove that $|O'(\tfrac12)|\ll W\cdot P$, where
$W=\sum_k e^{-\varepsilon^2\gamma_k^2}$ and $P=\pi(\kappa)$ is the
prime-counting function at the cutoff.
The numerical value at reference parameters is $O'(\tfrac12)=-2.475$,
approximately $11\%$ of $W\cdot P$, consistent with the cancellation
expected from the odd kernel $\sin(\gamma_k\log p)$.
The expected tools are: equidistribution of the sequence
$\{\gamma_k\log p\bmod 2\pi\}$ under appropriate averaging; a discrete
Weyl-type bound exploiting the independent behaviour of zero heights
and prime logarithms.
At the reference parameters, $B<0$ is numerically established
(Corollary~\ref{cor:curvature}); the remaining hypothesis of the
conditional remark following it is therefore $O'(\tfrac12)=0$.
A solution to this problem thus suffices on its own to discharge that
remark at the reference parameters and yield an unconditional strict
local-minimum statement for $O$ at $\sigma=\tfrac12$ at those
parameters.
\end{openproblem}

\begin{openproblem}[Asymptotic pointwise bias]
\label{op:bias-asymp}
Prove analytically that $\mathrm{Re}\,\Zt_p(\varepsilon)<0$ for every
prime $p\le\kappa$ in some explicit range $\varepsilon\in(0,\varepsilon_c)$,
together with the asymptotic ratio bound $r_p<1$ for every prime $p$.
At present both statements are verified numerically on the range
$p\le 53$, with all observed values of $r_p$ in $[0.75,0.80]$ and with
pointwise negativity established on the tested grid at
$\varepsilon=0.020$.
Closing this gap would upgrade Theorems~\ref{thm:ratio-signcross}
to proved statements and, in combination with an analytic bound on the
gamma contribution of Theorem~\ref{thm:gamma}, yield an analytic
version of the asymptotic bias statement for
$\mathrm{Re}\,\Zt_p(\varepsilon)$ as $\varepsilon\to 0$.
The expected tools are: explicit expansion of the pole residues of
$\xi$ at $s=0,1$ with the correct $\cosh$-evaluation of the test
function; a Riemann--Lebesgue-type bound on the digamma contribution
with explicit constant; a cancellation argument between the
$\varepsilon^{-1}$-coefficients of $\mathrm{Main}_p$ and
$\mathrm{Const}_p$.
The problem is self-contained within classical analytic number theory.
\end{openproblem}

\begin{openproblem}[Uniformity of positive curvature]
\label{op:uniformity}
Determine whether the inequality $B(\varepsilon,\kappa)<0$, and hence
$O''(\tfrac12)>0$, persists at parameter points
$(\kappa,\varepsilon,N)$ other than the reference values
$(53,0.05,100)$ tested here.
Corollary~\ref{cor:curvature} establishes the inequality at a single
point of the parameter space.
Whether it extends to a neighbourhood, to a larger range of cutoffs
$\kappa$, or to a different choice of smoothing $\varepsilon$ is a
distinct question: the direct curvature sum depends non-trivially on
the $\gamma_k^2$-weights and on the oscillatory phases $\gamma_k\log p$,
and finite parameter changes could in principle reverse the sign.
The expected tools are: a systematic numerical survey across a
parameter grid; an asymptotic expansion of $B(\varepsilon,\kappa)$ for
large $\kappa$; and, if possible, an analytic reduction via the
primal-to-dual identity $B=\sum_p(\log p)^2\,\mathrm{Re}\,Z_p$.
\end{openproblem}

\begin{openproblem}[Integrated-to-direct transfer]
\label{op:transfer}
Establish or disprove the implication
\[
\bigl(\,\mathrm{Re}\,\Zt_p(\varepsilon)<0 \text{ for all primes }
p\le\kappa\,\bigr)
\;\Longrightarrow\;
B(\varepsilon,\kappa)<0
\]
at fixed $(\varepsilon,\kappa)$, where
\[
B(\varepsilon,\kappa) \;:=\; \sum_{k,p} e^{-\varepsilon^2\gamma_k^2}\,
(\gamma_k\log p)^2\,\cos(\gamma_k\log p).
\]
All three quantities --- the pointwise sign of
$\mathrm{Re}\,\Zt_p$, the integrated bias $B_{\mathrm{int}}(\varepsilon)$,
and the direct curvature sum $B$ --- are numerically negative at the
reference parameters $(\kappa,\varepsilon,N)=(53,0.05,100)$
(Theorems~\ref{thm:ratio-signcross},~\ref{thm:intbias}, and
Corollary~\ref{cor:curvature}), but no analytical implication among
them is established.
A proof of the displayed implication would extend
Corollary~\ref{cor:curvature} from a reference-parameter statement to
a parametrised statement valid wherever the pointwise hypothesis can
be verified.
The expected tools are: the explicit prime-side representation of $B$,
the uniform positivity of the Fourier transform
$\widehat{h}_{p,\varepsilon}$, and a rearrangement argument
distributing the $\gamma_k^2$-weights across the pointwise
contributions.
A natural intermediate step would be the analogous implication with
$B_{\mathrm{int}}$ in place of $B$.
The problem is self-contained within classical analytic number theory.
\end{openproblem}

\begin{openproblem}[Lagarias--Weil bridge]
\label{op:weil-bridge}
The present result introduces a finite-cutoff curvature observable
for Gaussian-smoothed explicit-formula data; it should be compared
with Bombieri's finite truncations of Weil's quadratic
functional~\cite{Bombieri} and with the Li/Bombieri--Lagarias
positivity framework~\cite{Li,BombieriLagarias}, but it is not yet a
positivity theorem for the full Weil distribution.

We conjecture that the positive-curvature statement
$O''(\tfrac12)>0$, taken together with the stationarity condition
$O'(\tfrac12)=0$, implies Weil positivity on the relevant class of
smooth compactly-supported test functions in the formulation of
Lagarias~\cite{Lagarias}, and thereby implies the Riemann Hypothesis.
The conjecture is stated here as a named target rather than as a
proved implication; its resolution is beyond the methods of the
present paper and requires either a direct positivity argument on the
Weil distribution or an external route via Connes--Meyer-type operator
positivity.
Problems~\ref{op:stat}--\ref{op:transfer} are self-contained within
classical analytic number theory; the present problem is structurally
the principal outstanding question of the programme.
\end{openproblem}

% -------------------------------------------------------
\section{Non-Circularity}
\label{sec:nocirc}

We document explicitly that the results of this paper do not rely on
the objects they are directed toward.

\textbf{No RH as input.}
At no point is the Riemann Hypothesis assumed.
The ordinates $\gamma_k$ are used as numerically specified inputs, and
no property of their distribution requiring the Riemann Hypothesis is
assumed or used.

\textbf{No GUE, Montgomery, or Hilbert--P\'olya hypothesis.}
{\sloppy
The operator $\Tt$ and its $\sigma$-parametrised extension
$\Tt(\sigma)$ are constructed explicitly from Gaussian-smoothed prime
phasors in the previous papers of this series~\cite{Paper3,Paper4,Paper5}.
No random-matrix hypothesis, no pair-correlation conjecture, and no
Hilbert--P\'olya postulate enters the construction or the proofs.
\par}

\textbf{$\Tt$ is not a Hilbert--P\'olya operator.}
As established in~\cite{Paper5}, the eigenvalues of $\Tt$ decay as
$\mu_j\sim\CT/\gamma_{k(j)}\to 0$, which is incompatible with the
Hilbert--P\'olya requirement that eigenvalues grow like the ordinates
$\gamma_k$.
The paper does not attempt a Hilbert--P\'olya realisation.

\textbf{The Guinand--Weil formula is classical.}
Proposition~\ref{prop:gw} is a classical identity valid for every even,
sufficiently rapidly decaying test function; its derivation does not
require the truth of the Riemann Hypothesis.
Theorems~\ref{thm:main-neg},~\ref{thm:other}, and~\ref{thm:trunc},
together with the algebraic identity of Proposition~\ref{prop:curv-bias},
are proved unconditionally by direct computation.

\textbf{Numerical results are clearly marked.}
Theorems~\ref{thm:gamma},~\ref{thm:ratio-signcross},
and~\ref{thm:intbias}, and Corollary~\ref{cor:curvature}, rely on the
numerical evaluation of finite sums whose truncation error is
controlled by Theorem~\ref{thm:trunc}.
Each such statement is labelled [numerical] in its heading.

\textbf{Conditional statements are explicitly flagged.}
The conditional remark following Corollary~\ref{cor:curvature} is
stated conditionally on the stationarity hypothesis $O'(\tfrac12)=0$
(Problem~\ref{op:stat}) and adds no further unproved input beyond this
explicit assumption.
Problem~\ref{op:transfer} is named openly and is not used as a
hypothesis anywhere in \S\S\,\ref{sec:main}--\ref{sec:curvcor}.

\textbf{$\gamma_k$ as inputs, not conclusions.}
The ordinates appear as numerically specified inputs to the model.
The claim $\mathrm{Re}(\rho_k)=\tfrac12$ is not used anywhere; the
reference point $\sigma_0=\tfrac12$ is chosen as the object of study,
not assumed to be the unique zero location.

% -------------------------------------------------------
\section*{Acknowledgments}

The author thanks the project's structured review and verification
workflow for systematic error detection across draft versions, and
external pre-review for successive rounds of feedback that shaped the
final formulation of the results in this paper.

% -------------------------------------------------------
\section*{Call for Feedback}

This paper is part of an ongoing open research initiative toward the
Riemann Hypothesis.
All papers in the series are freely available on Zenodo under CC~BY~4.0,
and the accompanying verification code is hosted on GitHub.

The author welcomes critical feedback, corrections, and collaboration:
{\sloppy
Zenodo (\url{https://zenodo.org/search?q=Ulrich+Tehrani}) and GitHub
(\url{https://github.com/utehrani/analysislab-nt}).
\par}

Topics of particular interest include: analytical approaches to the
asymptotic pointwise bias of Problem~\ref{op:bias-asymp}; the analytic
stationarity bound of Problem~\ref{op:stat}; uniformity of the
curvature inequality across a larger parameter range
(Problem~\ref{op:uniformity}); the integrated-to-direct transfer of
Problem~\ref{op:transfer}; and, ultimately, the Lagarias--Weil bridge
of Problem~\ref{op:weil-bridge}.

We particularly welcome expert comment on the following question:
for a Gaussian-smoothed family $g_\varepsilon$ with
$g_\varepsilon\to\delta$ as $\varepsilon\to 0$, is it known whether
the sign of $W(g_\varepsilon * g_\varepsilon^\vee)$ at
finite~$\varepsilon$ carries arithmetic information, or is positivity
only meaningful in the distributional limit?

% -------------------------------------------------------
\begin{thebibliography}{99}

\bibitem{Paper1}
U.~Tehrani,
``A Curvature Decomposition of the Explicit Formula'',
Zenodo, \url{https://doi.org/10.5281/zenodo.19025598},
April 2026, v8.

\bibitem{Paper2}
U.~Tehrani,
``From Local Curvature to the Weil Functional'',
Zenodo, \url{https://doi.org/10.5281/zenodo.19106992},
April 2026, v5.

\bibitem{Paper3}
U.~Tehrani,
``A Finite-Cutoff Hilbert-Space Model
for Prime--Zero Energy Structure'',
Zenodo, \url{https://doi.org/10.5281/zenodo.19307989},
April 2026, v2.

\bibitem{Paper4}
U.~Tehrani,
``A Dual Operator for Prime--Zero Coupling
and a Conditional Proof of Energy Asymmetry'',
Zenodo, \url{https://doi.org/10.5281/zenodo.19364703},
April 2026, v4.

\bibitem{Paper5}
U.~Tehrani,
``Spectral Trace Formula and Smoothed Zero Sums:
A Prime--Zero Duality Framework'',
Zenodo, \url{https://doi.org/10.5281/zenodo.19508547},
April 2026, v1.1.

\bibitem{Code}
U.~Tehrani,
\emph{Verification scripts for the series},
GitHub, \url{https://github.com/utehrani/analysislab-nt}.

\bibitem{Lagarias}
J.C.~Lagarias,
``On a positivity property of the Riemann $\xi$-function'',
\emph{Acta Arithmetica} \textbf{89} (1999), 217--234.

\bibitem{Guinand}
A.P.~Guinand,
``A summation formula in the theory of prime numbers'',
\emph{Proc.\ London Math.\ Soc.}\ (2) \textbf{50} (1948), 107--119.

\bibitem{Weil}
A.~Weil,
``Sur les `formules explicites' de la th\'eorie des nombres
premiers'',
\emph{Comm.\ S\'em.\ Math.\ Univ.\ Lund, Suppl.}\ (1952), 252--265.

\bibitem{Hadamard}
J.~Hadamard,
``Sur la distribution des z\'eros de la fonction $\zeta(s)$
et ses cons\'equences arithm\'etiques'',
\emph{Bull.\ Soc.\ Math.\ France} \textbf{24} (1896), 199--220.

\bibitem{Edwards}
H.M.~Edwards,
\emph{Riemann's Zeta Function},
Academic Press, New York, 1974.

\bibitem{Connes}
A.~Connes,
``Trace Formula in Noncommutative Geometry and the Zeros of the
Riemann Zeta Function'',
\emph{Selecta Math.\ (N.S.)} \textbf{5}.1 (1999), 29--106.

\bibitem{BerryKeating}
M.V.~Berry and J.P.~Keating,
``The Riemann Zeros and Eigenvalue Asymptotics'',
\emph{SIAM Review} \textbf{41}.2 (1999), 236--266.

\bibitem{deBranges}
L.~de~Branges,
``The Riemann Hypothesis for Hilbert Spaces of Entire Functions'',
\emph{Bull.\ Amer.\ Math.\ Soc.\ (N.S.)} \textbf{15}.1 (1986),
1--17.

\bibitem{Burnol}
J.-F.~Burnol,
``On Fourier and Zeta(s)'',
\emph{Forum Math.} \textbf{16}.6 (2004), 789--840.

\bibitem{Bombieri}
E.~Bombieri,
``Remarks on Weil's Quadratic Functional in the Theory of Prime
Numbers, I'',
\emph{Rend.\ Lincei Mat.\ Appl.} \textbf{11}.3 (2000), 183--233.

\bibitem{Li}
X.-J.~Li,
``The Positivity of a Sequence of Numbers and the Riemann
Hypothesis'',
\emph{J.\ Number Theory} \textbf{65}.2 (1997), 325--333.

\bibitem{BombieriLagarias}
E.~Bombieri and J.C.~Lagarias,
``Complements to Li's Criterion for the Riemann Hypothesis'',
\emph{J.\ Number Theory} \textbf{77}.2 (1999), 274--287.

\bibitem{RubinsteinSarnak}
M.~Rubinstein and P.~Sarnak,
``Chebyshev's Bias'',
\emph{Experimental Math.} \textbf{3}.3 (1994), 173--197.

\end{thebibliography}

% -------------------------------------------------------
\bigskip\noindent\rule{\linewidth}{0.4pt}\par\smallskip
\noindent\small
\emph{Paper~6 v1.2 \textperiodcentered\ April~2026}\\
\emph{Literature positioning strengthened; eight references added;
operator-programme distinction, curvature comparison, and scope
clarification inserted; no changes to proofs or numerical values}\\
\emph{Pauca sed matura}

\end{document}
